\section{Purpose, assumptions, and how to read this document}
\label{sec:scope}

\subsection{The physical quantities of interest}

Let $\vR$ denote the vector collecting all nuclear coordinates.  The primary
objects studied here are obtained by first solving for the periodic restricted
Hartree--Fock (HF) density matrix $D_{\mathrm{HF}}(\vR)$ and then evaluating a
separate semilocal density functional approximation (DFA) on that same
density.  The resulting energy is
\begin{equation}
  E_{\mathrm{DC-DFA}}(\vR)
  = E_{\mathrm{DFA}}\!\left[D_{\mathrm{HF}}(\vR);\vR\right],
  \label{eq:scope-dc-energy}
\end{equation}
and its fixed-cell nuclear gradient is
\begin{equation}
  g_{Ax}
  = \frac{\dd E_{\mathrm{DC-DFA}}}{\dd R_{Ax}}.
  \label{eq:scope-gradient}
\end{equation}
Here, $A$ labels an atom and $x\in\{x,y,z\}$ is a Cartesian component.  The
public API returns the gradient; the force has the opposite sign,
\begin{equation}
  F_{Ax}=-g_{Ax}.
  \label{eq:scope-force-sign}
\end{equation}

Equation~\eqref{eq:scope-dc-energy} does not mean ``add a DFA correction to the
HF energy.''  The kinetic, electron--nuclear, Hartree, exchange--correlation,
and ion--ion terms of the target DFA energy are all re-evaluated on the HF
density.  The HF energy is retained only as the reference that generated the
source orbitals and as a diagnostic quantity.

\subsection{Why this differs from an ordinary KS gradient}

In a self-consistent KS-DFA calculation, the same DFA energy is stationary with
respect to orbital rotations $\kappa$.  In DC-DFA, the functional that defines
the orbitals differs from the functional used to evaluate the energy, so that
\begin{equation}
  \left.
  \frac{\partial E_{\mathrm{DC-DFA}}}{\partial\kappa_{ai}^{\vk}}
  \right|_{D_{\mathrm{HF}}}
  \ne 0.
  \label{eq:scope-nonstationary}
\end{equation}
The change of the HF orbitals under nuclear motion therefore cannot be
discarded.  Rather than solve a separate CPHF equation for every nuclear
displacement, one solves a single adjoint equation for the HF orbital Hessian,
the Z-vector equation, and contracts that solution with all perturbations.

\subsection{Active scope}

This edition treats the following as one complete theoretical and numerical
scope:
\begin{itemize}
  \item restricted closed-shell RHF/KRHF source densities,
  \item three-dimensional periodic boundary conditions,
  \item GTH pseudopotentials and atom-centered Gaussian AOs,
  \item Gamma-point and general sampled k-point meshes,
  \item PBE, SCAN, and r2SCAN target DFAs,
  \item FFTDF, GDF, and RSDF two-electron representations,
  \item GPU target multigrid XC or J+XC selected independently of source DF,
  \item dense-reference and production matrix-free Z-vectors,
  \item CPU and GPU implementations, and
  \item energies and nuclear analytic gradients at fixed lattice vectors.
\end{itemize}

Stress and cell derivatives, multigrid pseudopotentials, the larger-angular-momentum
completion matrix, and open-shell/UHF response are excluded from the active
scope and preserved in a separate archive.  Accordingly, ``analytic gradient''
always denotes $\partial E/\partial R_{Ax}$ at fixed cell in this document.

\subsection{Prerequisites and organization}

The reader is assumed to know AO-based HF and KS-DFT, density matrices,
Gaussian integrals, and elementary reciprocal-lattice theory.
Section~\ref{sec:notation} fixes the periodic notation and units, and
Sec.~\ref{sec:fixed-density-energy} defines the target energy evaluated on a
fixed HF density.  Sections~\ref{sec:orbital-response} and
\ref{sec:dc-lagrangian} derive the nonvariational response and the total
gradient.  Subsequent sections develop the explicit derivatives for GTH,
FFTDF, GDF, MGGA, and RSDF, followed by target-multigrid ownership,
matrix-free and bounded-memory response algorithms, a worked workflow, and
numerical validation procedures.

The word ``exact'' must be interpreted at two distinct levels.  First, within
a chosen finite basis, mesh, auxiliary basis, and k-point representation, the
analytic derivative must agree with the full-energy finite difference of the
same representation.  Second, the finite representation itself must converge
with cutoff, basis, and k-mesh.  The first criterion measures implementation
correctness; the second measures physical convergence.
